% ---------------------------------------------------------
\subsection{Axiom of Completeness}
% ---------------------------------------------------------


\begin{remark*}[Why completeness is needed]
The ordered field axioms alone do not prevent ``holes'' in the number line.
The rationals $\mathbb{Q}$ satisfy all field and order axioms, yet fail
completeness. Consider the set
\[
S = \{x \in \mathbb{Q} : x^2 < 2\}.
\]
This set is nonempty (e.g.\ $1 \in S$) and bounded above in $\mathbb{Q}$
(e.g.\ $2$ is an upper bound). Yet $\sup S$ does not exist as a rational
number: the candidate $\sqrt{2}$ is irrational. The set $S$ has no least
upper bound in $\mathbb{Q}$ --- a hole sits exactly where $\sqrt{2}$ should be.

The Completeness Axiom asserts that $\mathbb{R}$ has no such holes:
every nonempty bounded-above set has a supremum \emph{inside} $\mathbb{R}$.
This single axiom is what distinguishes $\mathbb{R}$ from $\mathbb{Q}$,
and it underlies every major theorem in analysis.
\end{remark*}


\begin{tcolorbox}[colback=propbox, colframe=propborder, arc=2pt,
  left=6pt, right=6pt, top=4pt, bottom=4pt,
  title={\small\textbf{Definition (Least Upper Bound Property)}},
  fonttitle=\small\bfseries]
\begin{definition}
\label{def:least-upper-bound-property}
Let $S$ be a partially ordered set. The set $S$ is said to have the
\emph{Least Upper Bound Property} if every nonempty subset of $S$ that is bounded
above in $S$ has a least upper bound in $S$.
\end{definition}
\end{tcolorbox}

\begin{remark*}[Standard quantified statement]
\[
\begin{aligned}
&S \text{ has the Least Upper Bound Property}
\\
&\qquad\Longleftrightarrow\qquad
\forall A \subseteq S\;
\Bigl[
(A \neq \varnothing \land \exists u \in S\;\forall x \in A\;(x \le u))
\\
&\qquad\qquad\rightarrow
\exists s \in S\;
\Bigl(
(\forall x \in A\;(x \le s))
\land
(\forall u \in S\;((\forall x \in A\;(x \le u)) \rightarrow s \le u))
\Bigr)
\Bigr].
\end{aligned}
\]
\end{remark*}

\begin{remark*}[Predicate statement]
\[
LeastUpperBoundProperty(S).
\]
Equivalently,
\[
\begin{aligned}
&\forall A\;
\Bigl[
(Subset(A,S) \land NonemptySet(A) \land BoundedAbove(A))
\\
&\qquad\rightarrow
\exists s \in S\; Supremum(s,A)
\Bigr].
\end{aligned}
\]
\end{remark*}

\begin{remark*}[Negated quantified statement]
\[
\begin{aligned}
&S \text{ does not have the Least Upper Bound Property}
\\
&\qquad\Longleftrightarrow\qquad
\exists A \subseteq S\;
\Bigl[
(A \neq \varnothing \land \exists u \in S\;\forall x \in A\;(x \le u))
\\
&\qquad\qquad\land
\forall s \in S\;
\Bigl(
(\exists x \in A\;(x > s))
\lor
(\exists v \in S\;((\forall x \in A\;(x \le v)) \land v < s))
\Bigr)
\Bigr].
\end{aligned}
\]
This means there exists a nonempty subset of $S$ that is bounded above in $S$ but has no
least upper bound in $S$. The failure is not in the subset alone, but in the ambient ordered
system $S$.
\end{remark*}

\begin{remark*}[Interpretation]
The Least Upper Bound Property says that whenever a nonempty subset of $S$ is bounded
above, it has a smallest upper bound in $S$.

The important point is that the issue is not in the subset alone. A bounded set can fail to
have a least upper bound only because $S$ is missing the element that should serve as that
bound.
\end{remark*}



\begin{tcolorbox}[colback=axiombox, colframe=axiomborder, arc=2pt,
  left=6pt, right=6pt, top=4pt, bottom=4pt,
  title={\small\textbf{Axiom (Completeness of $\mathbb{R}$)}},
  fonttitle=\small\bfseries]
\begin{axiom}
\label{ax:completeness-of-reals}
The ordered field $\mathbb{R}$ has the Least Upper Bound Property.
\end{axiom}
\end{tcolorbox}

\begin{remark*}[Standard quantified statement]
\[
\begin{aligned}
&\mathbb{R} \text{ has the Least Upper Bound Property}
\\
&\qquad\Longleftrightarrow\qquad
\forall A \subseteq \mathbb{R}\;
\Bigl[
(A \neq \varnothing \land \exists u \in \mathbb{R}\;\forall x \in A\;(x \le u))
\\
&\qquad\qquad\rightarrow
\exists s \in \mathbb{R}\;
\Bigl(
(\forall x \in A\;(x \le s))
\land
(\forall u \in \mathbb{R}\;((\forall x \in A\;(x \le u)) \rightarrow s \le u))
\Bigr)
\Bigr].
\end{aligned}
\]
\end{remark*}

\begin{remark*}[Predicate statement]
\[
LeastUpperBoundProperty(\mathbb{R}).
\]
Equivalently,
\[
\begin{aligned}
&\forall A\;
\Bigl[
(Subset(A,\mathbb{R}) \land NonemptySet(A) \land BoundedAbove(A))
\\
&\qquad\rightarrow
\exists s \in \mathbb{R}\; Supremum(s,A)
\Bigr].
\end{aligned}
\]
\end{remark*}

\begin{remark*}[Negated quantified statement]
\[
\begin{aligned}
&\mathbb{R} \text{ does not have the Least Upper Bound Property}
\\
&\qquad\Longleftrightarrow\qquad
\exists A \subseteq \mathbb{R}\;
\Bigl[
(A \neq \varnothing \land \exists u \in \mathbb{R}\;\forall x \in A\;(x \le u))
\\
&\qquad\qquad\land
\forall s \in \mathbb{R}\;
\Bigl(
(\exists x \in A\;(x > s))
\lor
(\exists v \in \mathbb{R}\;((\forall x \in A\;(x \le v)) \land v < s))
\Bigr)
\Bigr].
\end{aligned}
\]
This is the hypothetical failure of completeness: it asserts that there exists a nonempty subset
of $\mathbb{R}$ that is bounded above but has no least upper bound in $\mathbb{R}$. The failure
would lie not in the subset alone, but in the ambient ordered field.

A standard example is the set
\[
A=\{q \in \mathbb{Q} : q^2 < 2\}.
\]
This set is nonempty and bounded above in $\mathbb{Q}$, but it has no supremum in $\mathbb{Q}$; the missing least upper bound is $\sqrt{2}$, so the failure lies in the field $\mathbb{Q}$, not in the set itself.
\end{remark*}

\begin{remark*}[Equivalent formulations]
Over the ordered-field axioms, completeness is equivalent to each of the
following:
\begin{itemize}
  \item Every nonempty set bounded below has an infimum.
  \item Every Cauchy sequence in $\mathbb{R}$ converges.
  \item \emph{(Nested Interval Property)} Every nested sequence of nonempty
        closed bounded intervals has nonempty intersection.
\end{itemize}
\end{remark*}
